\section{Implications for Evaluation Validity}
A high-scoring evaluation-awareness probe is evidence about a measurement instrument, not automatically evidence that a model represents evaluation purpose independently of benchmark structure. Equal-N training, matched payload blocks, separate AUC and balanced accuracy, and independent endpoint validation should be treated as validity checks. The practical output is a claim ladder: format shortcut, purpose decodability, equal-N survival, geometry, and only then causal transport.

The mechanistic result is the paper's main contribution. Naive probes can exploit prompt format; equal-N decorrelation shows that pooled gains are not only a row-count artifact; and the matched $2\times2$ design decomposes purpose representation into shared purpose ($A$), format ($B$), and purpose$\times$format interaction ($C$). Geometry selected L5 as shared and L16 as interaction-heavy before intervention. At L16, correctly signed $C$ improved transport toward the model's natural matched purpose counterfactual, while wrong-sign, shuffled, norm-matched, and fixed-random controls did not. This is an internal representation result under controlled prompt-provided context, not evidence that the model spontaneously recognizes evaluation or can be behaviorally steered.
